\section{Limitations}\label{sec:limits}
\begin{itemize}
\item \textbf{Simulation only, one dwelling, one typical year.} The plant is a reduced-order
RC model. It has not been calibrated against measurements or cross-validated against a
detailed whole-building simulation. An EnergyPlus model with the same geometry,
constructions, gains and EPW file is released with the code
(\texttt{validation/energyplus}), together with a script that compares free-floating
temperatures and ideal cooling loads, but that comparison had not been run when this
paper was written. The
absolute energy figures should therefore be read as indicative; the comparisons
between controllers, made on identical inputs, are more robust. Cross-validation
against EnergyPlus is the most important next step.
\item \textbf{Assumed inputs.} Envelope properties, gains, occupancy and equipment
performance are assumed. A 40-sample variation of eleven of them leaves the controller
ranking unchanged but changes the size of the savings (Section~\ref{sec:sens}).
\item \textbf{Tuning on the evaluation year.} The E-MPC weights and back-off, the
Hybrid gains and the DQN hyper-parameters were chosen during development by running on
the same typical year that is used for evaluation. The DQN's validation weeks are also
drawn from that year. The reported margins may therefore be slightly optimistic for an
unseen year; the heatwave case is a partial out-of-sample test.
\item \textbf{Assumed equipment curves.} The VRF rated COP, the part-load curve and the
IEC effectiveness are generic values, not manufacturer data. The released configuration
file allows them to be replaced.
\item \textbf{Representative outdoor PM$_{2.5}$.} The outdoor series is a monthly--diurnal
profile, not measured data. The IAQ results are therefore comparisons between
architectures, not exposure estimates. Hourly CPCB monitoring data for Pune should
replace it.
\item \textbf{Single-volume IAQ.} CO$_2$, PM$_{2.5}$ and moisture are modelled as one
well-mixed volume. Bedrooms with closed doors at night would see higher CO$_2$ than the
whole-dwelling value reported here.
\item \textbf{Dry-bulb comfort.} The comfort band is a static 23--26\,$^\circ$C. Indoor
humidity is reported separately rather than combined into PMV or an adaptive (IMAC)
index.
\item \textbf{Deterministic schedules.} Occupancy, internal gains and cooking follow fixed
weekday and weekend profiles, and occupants cannot override the controller.
\item \textbf{DQN scope.} One value-based algorithm with a discrete action set, trained
in simulation on perturbed copies of the same typical year. Policy-gradient and
continuous-action methods, and training on multiple measured years, were not tested.
\item \textbf{No heating.} The dwelling has no heating, so winter cold is outside the
controllers' authority. It is excluded from the discomfort metric through the
free-floating reference.
\item \textbf{Economics.} The main three-tier ToD tariff is hypothetical, not a published
MSEDCL schedule. The actual MSEDCL residential structure, a flat tariff and a steeper
tariff are tested as sensitivities. One national grid emission factor is used
(CEA, 0.727\,kgCO$_2$/kWh for FY2023--24). Demand charges and tariff revisions are not modelled.
\end{itemize}
